% Section 10.4, from lectures/L10/appendix/c_bringup.md.

\section{Bringing Up the Bench}
\label{c10:sec:exercises}\label{c10:app:c}

\begin{aside}
\textbf{How to work these.} This is the bench work that ends the book: the DE0-CV, the Nano, a level
shifter and a USB-serial adapter, wired into one working node. The hardware is listed in the
preface's \emph{Setting up}. Read \secref{c10:sec:bench} first: it explains the two links, the level
shifting, and why the bring-up ladder climbs one rung at a time. \code{app::EchoNode}
(Exercise~\ref{c10:ex:b1}) should already be written and passing its host test.

\textbf{How to check your work.} The bench is the check: every rung of the ladder has an observable
outcome, and the first one that misbehaves names the layer at fault. The answers to the parts that
are not bench work are in \secref{sol:c10}.
\end{aside}

% ------------------------------------------------------------------------------------------
\exercise{Assemble the bench}{Bench}
\label{c10:ex:c1}
Wire it with everything \textbf{powered off}, and have the wiring checked before applying power. A
single 5\,V line into a 3.3\,V FPGA pin can damage it.

\textbf{The control plane (SPI), through the level shifter.} The Nano's SPI pins are fixed; the
DE0-CV signals are the GPIO-header pins assigned in \file{uart_board.vhd}.

\begin{narrowtable}{@{}llll@{}}
  \thead{Signal} & \thead{Nano} & \thead{Level shifter} & \thead{DE0-CV} \\
  \midrule
  \code{SCK} & D13 (PB5) & HV to LV & \code{sclk} \\
  \code{MOSI} & D11 (PB3) & HV to LV & \code{mosi} \\
  \code{MISO} & D12 (PB4) & LV to HV & \code{miso} \\
  \code{SS}, the chip select & D10 (PB2) & HV to LV & \code{ss} \\
  5\,V reference & 5V & HV & --- \\
  3.3\,V reference & --- & LV & 3.3\,V \\
  GND & GND & GND (both sides) & GND \\
\end{narrowtable}

\textbf{The data plane (UART), 3.3\,V direct}, with no shifter, since the adapter is set to 3.3\,V
logic. The \code{tx}/\code{rx} crossover is the usual one.

\begin{narrowtable}{@{}lll@{}}
  \thead{Signal} & \thead{DE0-CV} & \thead{USB-serial adapter} \\
  \midrule
  The peripheral's \code{tx} & \code{tx} & RX \\
  The peripheral's \code{rx} & \code{rx} & TX \\
  GND & GND & GND \\
\end{narrowtable}

The adapter's USB goes to the PC, where a terminal at 115200~8N1 is the far end of the data plane.

% ------------------------------------------------------------------------------------------
\exercise{Program the FPGA and flash the Nano}{Bench}
\label{c10:ex:c2}
Two chips, two toolchains, both of which you have already used. For the \textbf{DE0-CV},
re-synthesize \file{uart_board.vhd}, the provided Quartus top level wrapping your \code{uart_top},
and program the FPGA over USB-Blaster exactly as in \chapref{9}; confirm that the design fits and
that the pin assignments match the wiring table above. For the \textbf{Nano}, flash the ATmega
firmware:

\begin{shell}
make -C fw/avr flash        # avr-gcc + avrdude; see fw/README.md
\end{shell}

\noindent With both programmed and the bench wired, you are ready to climb the ladder.

% ------------------------------------------------------------------------------------------
\exercise{Climb the bring-up ladder}{Bench}
\label{c10:ex:c3}
Climb one rung at a time and \textbf{stop at the first that misbehaves}; that rung names the layer at
fault (\secref{c10:sec:ladder} explains why).

\xpart{a}\parttitle{Data-plane pin loopback.} This is the loopback you already ran on the board alone
in \chapref{9}. In \file{uart_board.vhd}, temporarily wire the board's UART receive pin straight back
out to its transmit pin, bypassing \code{uart_top} altogether, and re-synthesize. Type in the PC
terminal; every character should come straight back. None of your logic is in this path, and that is
the point: it proves the USB-serial adapter, its 3.3\,V logic levels, the two data-plane wires, and the FPGA pin
assignment, and it proves nothing else, because nothing else is in the circuit yet. Note what it
cannot prove: the terminal's baud and frame settings. The adapter sends and receives at the same
settings, so a loopback agrees with itself at any rate; they are first on trial in rung d. Fix any problem here before going further, since every later
rung depends on this one.

\xpart{b}\parttitle{The control plane.} Restore the wrapper. Flash a bring-up \code{main()} that does
nothing but write \code{BAUD_DIV} over SPI and read it straight back, comparing the two values. A
matching read-back is the first evidence that the whole control plane works: the level shifter on all
four SPI lines, \code{AvrSpi}'s register setup, the provided \code{spi_slave} and
\code{spi_reg_bridge}, and your register bank's read path. This rung has to come before any rung
that uses the peripheral, because until \code{BAUD_DIV} is written the divider reads zero and
\code{uart_top} substitutes 1, so a byte written to \code{TX_DATA} would leave at 50\,MHz / 16 =
3.125\,Mbaud --- transmitted, but at a rate nothing on the bench can decode --- and the only route to
a usable \code{BAUD_DIV} is SPI. A crossed, unshifted or floating SPI line surfaces here and nowhere
earlier.

\xpart{c}\parttitle{The peripheral, looped back on itself.} Wire the peripheral's own \code{tx} to
its own \code{rx} in the wrapper, exactly as \code{uart_top_tb} does in simulation, and
re-synthesize. Have the bring-up \code{main()} configure the peripheral, write a byte to
\code{TX_DATA}, then poll \code{STATUS} and read \code{RX_DATA} back. The same byte coming back
proves the entire datapath on real silicon: \code{baud_gen}, \code{uart_tx}, \code{uart_rx}, both
FIFOs and the register bank, running at an actual 50\,MHz rather than a simulated one. The PC
terminal is not involved at all, so a failure here is the FPGA and not the link.

\xpart{d}\parttitle{The real data plane.} Undo the loopback and wire \code{tx} and \code{rx} to the
USB-serial adapter. Repeat the send and the receive from rung c, but against the terminal this time:
a byte the driver writes should appear in the terminal, and a character you type should come back
through \code{RX_DATA}. This adds the one thing rung c could not, which is the peripheral talking to a
\emph{different} device, with two independently generated baud rates that now have to agree with
each other.

\xpart{e}\parttitle{\code{app::EchoNode}.} Flash the echo application you wrote and host-tested in
Exercise~\ref{c10:ex:b1}: \code{main()} constructs a \code{Uart} over \code{AvrSpi},
\code{configure()}s it for 115200~8N1, and calls \code{node.run(stop)}. What you type in the terminal
is received by the driver over SPI and echoed back out the peripheral's \code{tx}, one byte crossing
every layer. Because the class was already proven on the host, anything that fails at this rung is
the wiring or the link, not the echo logic.

% ------------------------------------------------------------------------------------------
\exercise{Trace a byte}{Reasoning}
\label{c10:ex:c4}
The closing exercise, and the book's summary. Take a single character typed in the terminal and
trace it through \textbf{every} layer, out and back, naming each module and the chapter that built
it.

It travels in on the data plane, from the PC terminal to the USB-serial adapter to the DE0-CV's
\code{rx} pin, through \code{sync} into \code{uart_rx} (both \chapref{3}), then the RX FIFO
(\code{fifo}, \chapref{4}) inside \code{uart_regs} (\chapref{5}). It then crosses the control plane,
out through \code{spi_reg_bridge} and \code{spi_slave} (provided) to \code{AvrSpi} (\chapref{9}), up
through \code{readReg} and \code{Uart} (\chapref{7}), and into \code{EchoNode} (\chapref{10}).
Finally it goes back out, from \code{writeBlocking()} (\chapref{8}) and \code{Uart::write()} to
\code{writeReg()} to \code{AvrSpi}, back across \code{spi_slave} and \code{spi_reg_bridge} into
\code{uart_regs} and the TX FIFO, through the TX feeder in \code{uart_top} (\chapref{5}) and
\code{uart_tx} (\chapref{2}) to the DE0-CV's \code{tx} pin, and back via the adapter to the terminal.

Name each place the byte changes representation (a UART frame on the wire, a register field over
SPI, a \code{uint8_t} in the driver), and you have described the whole FPGA-meets-MCU stack the book
built.
